\section{多重积分与变量替换}
\label{sec:02-Calculus/06-Multivariable-Calculus/07}

从标准正态里抽一个样本 \(\mathbf z\)，送进一个可逆网络 \(T\)，得到 \(\mathbf x=T(\mathbf z)\)。新样本的密度是多少？

一个常见的直觉是``密度跟着搬家''：\(\mathbf z\) 处密度是多少，像点 \(\mathbf x\) 处就是多少。这个直觉错在它忘了密度是``每单位体积的概率''。如果 \(T\) 在 \(\mathbf z\) 附近把小方块撑成了三倍大的平行六面体，同样多的概率摊在三倍的体积上，密度就得除以三。要把这句话变成公式，需要一个能回答``局部体积被放大了多少倍''的量。本节给出这个量，而它早在两节之前就出现过了。

\subsection{累积的骨架：局部量乘小体积，再求和}

把区域 \(D\subset\R^2\) 切成许多小块，在第 \(i\) 块里取一点，作和
\[
\sum_i f(x_i^*,y_i^*)\,\Delta A_i .
\]
细分加密时若这个和收敛到与切法、取点无关的极限，就记作 \(\iint_D f\,dA\)。三重积分同理，把 \(\Delta A\) 换成 \(\Delta V\)。

整个定义只在说一件事：\emph{局部强度乘以小块尺寸，然后累加}。这决定了它的读法。\(f\) 是面密度时积分给总质量；\(f\) 是概率密度时给落在 \(D\) 里的概率；\(f\equiv1\) 时给面积本身。量纲检查因此格外好用——被积函数的单位乘上 \(dA\) 的单位，必须等于你想要的那个量的单位。

计算上通常化成迭代积分。对
\[
D=\set{(x,y):a\le x\le b,\;g_1(x)\le y\le g_2(x)},
\]
在合适的可积条件下
\[
\iint_D f\,dA=\int_a^b\!\!\left[\int_{g_1(x)}^{g_2(x)}f(x,y)\,dy\right]dx,
\]
内层沿一条竖直切片累积，外层把切片相加。Fubini 定理允许交换顺序，但交换的是\emph{切片方式}，不是两个字母的位置：先对 \(x\) 积分时，边界必须按水平切片重新描述。三角形 \(\set{0\le x\le1,\;0\le y\le1-x}\) 换个方向就得写成 \(\set{0\le y\le1,\;0\le x\le1-y}\)。忘了改边界是这一步唯一常见的错误，也是唯一值得记的判断点；至于内层外层怎么算，那是单变量的事。Fubini 何时合法、非负性与绝对可积性各自起什么作用，见\hyperref[sec:12-Real-Analysis-Support/06-Product-Measures-and-Fubini/02]{``Fubini 与 Tonelli：二重积分何时能拆成两次''}。

\subsection{换坐标必须补上体积伸缩}

设 \((x,y)=T(u,v)\) 是可微的坐标变换。\(uv\) 平面上以 \((u,v)\) 为角的小矩形 \(du\,dv\)，经过局部线性化被送成一个平行四边形，它的两条边是 \(J_T\) 的两列各乘上 \(du\)、\(dv\)。平行四边形的面积等于这两列所张成的面积，也就是行列式的绝对值乘以原面积。

\begin{center}
\begin{tikzpicture}[x=0.95cm,y=0.95cm,font=\small,>=stealth]
  \draw[black!18,step=0.5] (0,0) grid (2,2);
  \draw[->,black!45] (-0.2,0) -- (2.4,0) node[right] {\(u\)};
  \draw[->,black!45] (0,-0.2) -- (0,2.5) node[above] {\(v\)};
  \fill[blue!16] (1,1) rectangle (1.5,1.5);
  \draw[blue!70!black,thick] (1,1) rectangle (1.5,1.5);
  \node[below,black!65] at (1.25,0.9) {面积 \(du\,dv\)};

  \draw[->,thick,black!65] (2.7,1.6) -- (3.9,1.6);
  \node[above,black!65] at (3.3,1.65) {\(T\)};

  \begin{scope}[shift={(4.3,0)}]
    \draw[->,black!45] (-0.2,0) -- (4.4,0) node[right] {\(x\)};
    \draw[->,black!45] (0,-0.2) -- (0,3.1) node[above] {\(y\)};
    \draw[black!18] (0,0) -- (1.2,2.2);
    \draw[black!18] (0.7,0.15) -- (1.9,2.35);
    \draw[black!18] (1.4,0.3) -- (2.6,2.5);
    \draw[black!18] (2.1,0.45) -- (3.3,2.65);
    \draw[black!18] (2.8,0.6) -- (4.0,2.8);
    \draw[black!18] (0,0) -- (2.8,0.6);
    \draw[black!18] (0.3,0.55) -- (3.1,1.15);
    \draw[black!18] (0.6,1.1) -- (3.4,1.7);
    \draw[black!18] (0.9,1.65) -- (3.7,2.25);
    \draw[black!18] (1.2,2.2) -- (4.0,2.8);
    \fill[blue!16] (2.0,1.4) -- (2.7,1.55) -- (3.0,2.10) -- (2.3,1.95) -- cycle;
    \draw[blue!70!black,thick] (2.0,1.4) -- (2.7,1.55) -- (3.0,2.10) -- (2.3,1.95) -- cycle;
    \node[below right,black!65] at (2.75,1.42) {面积 \(\abs{\det J_T}\,du\,dv\)};
  \end{scope}
\end{tikzpicture}

\emph{同一张网格在两侧对应。直线仍是直线、平行仍然平行，这是线性近似的功劳；变的是每个小格的面积，倍率由 \(\abs{\det J_T}\) 在该点的取值给出，并且随点变化。}
\end{center}

\begin{theorem}[变量替换]
若 \(T\) 在开集上连续可微、\(\det J_T\ne0\) 且在所涉区域上（除去一个零面积的例外集）一一对应，则
\[
\iint_D f(x,y)\,dx\,dy
=\iint_{T^{-1}(D)}f(T(u,v))\,\abs{\det J_T(u,v)}\,du\,dv .
\]
\end{theorem}

绝对值保证面积非负；行列式本身的符号只记录定向有没有翻转，那是曲面积分和微分形式关心的事。漏掉这个因子，等于默认新坐标的小格和旧坐标的小格一样大——恰恰是坐标变换最不可能保持的性质。

极坐标是最省事的检验。\(x=r\cos\theta\)、\(y=r\sin\theta\) 的 Jacobian 已经在\hyperref[sec:02-Calculus/06-Multivariable-Calculus/04]{``Jacobian 与多元链式法则''}一节按列读过，取行列式得 \(r\)，于是 \(dA=r\,dr\,d\theta\)。几何上一目了然：小扇环的径向厚度是 \(dr\)，切向长度约为 \(r\,d\theta\)，离原点越远同样的角度扫过的弧越长。半径 \(R\) 的圆盘面积随之为 \(\int_0^{2\pi}\!\int_0^R r\,dr\,d\theta=\pi R^2\)；丢了 \(r\)，答案会退化成与 \(R\) 一次方成正比的荒唐结果。柱坐标的体积元是 \(r\,dr\,d\theta\,dz\)，球坐标在常用约定下是 \(\rho^2\sin\phi\,d\rho\,d\phi\,d\theta\)。这些不值得背——写出坐标映射、求一次行列式即可，而且这样做还会逼你先把角度的约定说清楚。

\subsection{概率密度变换是同一条定理}

现在回到开头。设随机向量 \(\mathbf Z\) 有密度 \(p_Z\)，\(\mathbf X=T(\mathbf Z)\) 且 \(T\) 可逆可微。对任意区域 \(A\)，
\[
P(\mathbf X\in A)=P\bigl(\mathbf Z\in T^{-1}(A)\bigr)
=\int_{T^{-1}(A)}p_Z(\mathbf z)\,d\mathbf z .
\]
在右边作代换 \(\mathbf z=T^{-1}(\mathbf x)\)，按上面的定理配上体积因子，得到
\[
p_X(\mathbf x)=p_Z\bigl(T^{-1}(\mathbf x)\bigr)\,\abs{\det J_{T^{-1}}(\mathbf x)}
=\frac{p_Z(\mathbf z)}{\abs{\det J_T(\mathbf z)}},\qquad \mathbf z=T^{-1}(\mathbf x).
\]
所谓``密度要除以体积伸缩''，不是概率论额外规定的一条法则，它就是变量替换定理。同一条式子在\hyperref[sec:05-Probability/05-Joint-Distributions/05]{``多维变量变换与 Jacobian''}一节里从概率一侧再讲一遍。

取对数就是流模型的训练目标：
\[
\log p_X(\mathbf x)=\log p_Z(\mathbf z)-\log\abs{\det J_T(\mathbf z)} .
\]
这解释了这一族模型的整个设计压力所在。要让对数似然可算，就必须能便宜地求出 \(\log\abs{\det J_T}\)。一个一般的 \(d\times d\) 行列式要 \(O(d^3)\)，在 \(d\) 上千时每步都算一次是不现实的。于是各种耦合层、自回归结构被设计出来，它们的共同点是把 Jacobian 逼成三角矩阵——三角阵的行列式是对角元之积，代价降到 \(O(d)\)。表面上是网络结构的技巧，底下是这一节的定理在收费。

\subsection{定理不管的地方}

变换需要在所用区域上足够光滑，并且除掉一个可以忽略的例外集之外一一对应。若同一个 \((x,y)\) 被多个 \((u,v)\) 覆盖，积分会重复计数——极坐标的 \(\theta\) 加上 \(2\pi\) 回到同一点，正是为什么必须把角度限制在一个周期内。

退化点则要看它有多大。极坐标在原点处 \(\det J_T=r=0\) 且角度不唯一，但那只是一个点，面积为零，对积分没有影响。若退化发生在一整条曲线甚至一片区域上，情况就不同了：那里的映射把二维压成了更低维，逆映射不存在，必须把这部分单独挖出来处理。同样的警示适用于流模型——网络的 Jacobian 行列式在训练中接近零，意味着映射正在把某个方向压扁，对数似然会趋于正无穷而数值上先崩掉。

累积的第一步到此完成：我们已经能把标量在区域上加总，并且在换坐标时把账算对。但很多量并不摊在一片区域上，而是沿一条路径分布——沿航线做的功、沿路径累积的代价。下一节把累积搬到曲线上，并会碰到一个新问题：结果到底依不依赖走的是哪条路。
